%Root main.tex
%---------------------------------------------------------------------------
%付録E　付録
%---------------------------------------------------------------------------

\chapter{MMCコントローラ書き込み方法}
\thispagestyle{fancy}
\label{chap:}

\section{ロジックアナライザの設定}

(1)~~観測する信号があるモジュールを右クリックし，直下にicon，ilaのIPブロックを作成する.

\begin{figure}[H]
	\centering
	\includegraphics[width=110mm]{gaze/app_E/IPbrock}
	\caption{chipscope}
%	\label{fig:　}
\end{figure}

\begin{figure}[H]
	\centering
	\includegraphics[width=110mm]{gaze/app_E/new_source_wizard}
	\caption{chipscope}
%	\label{fig:　}
\end{figure}

(2)~~(1)のモジュールのソースコードにiconとilaのcomponent及びportmapを記述する．\\
*iconは図のように記述する．
*ilaのclkは内部信号の取得周期となる．
*ilaのTRIGに任意の信号を入力する

\begin{figure}[H]
	\centering
	\includegraphics[width=110mm]{gaze/app_E/ILA_1}
	\caption{chipscope}
%	\label{fig:　}
\end{figure}

(3)~~ILAをダブルクリックして開き，Number of trigger portsの数をTRIGの数と一致させる．Sapmple Data Depthは何周期分のデータを取得するか決定する．

\begin{figure}[H]
	\centering
	\includegraphics[width=110mm]{gaze/app_E/ILA_2}
	\caption{chipscope}
%	\label{fig:　}
\end{figure}

(4)~~各TRIGのビット幅を(2)で宣言した幅に合わせる．

\begin{figure}[H]
	\centering
	\includegraphics[width=110mm]{gaze/app_E/port_map}
	\caption{chipscope}
%	\label{fig:　}
\end{figure}

(5)~~その後，Generateをクリックする．


\newpage
\section{書き込み方法}
(1)~~コントローラに接続されている24V電源以外の物を全て外す.\\
*全て接続したまま書き込みも可能.

(2)~~ISEをインストールしたPCとコントローラをJTAGで接続する.

(3)~~コントローラの24V電源を入れる.

(4)~~ISEを起動し,File$\rightarrow$Open Projectからプログラムを開く.

\begin{figure}[H]
	\centering
	\includegraphics[width=110mm]{gaze/app_E/4}
	\caption{ISE File画面}
%	\label{fig:　}
\end{figure}

(5)~~左側タブ内の以下の項目について緑色のチェックマーク(または注意△！マーク)になっていることを確認する.\\
Synthesize –XST\\
Check Syntax\\
Implement Design\\
Generate Programming File\\
もしも,チェックされていない場合は,それぞれ右クリックし,Runをクリックする.\\

\begin{figure}[H]
	\centering
	\includegraphics[width=110mm]{gaze/app_E/5_1}
	\caption{ISE  Process画面1}
%	\label{fig:　}
\end{figure}

\newpage

\begin{figure}[H]
	\centering
	\includegraphics[width=110mm]{gaze/app_E/5_2}
	\caption{ISE  Process画面2}
%	\label{fig:　}
\end{figure}

(6)~~Generate Target PROM/ACE Fileをダブルクリックする.

(7)~~WarningダイアログはOKをクリックする.

\begin{figure}[H]
	\centering
	\includegraphics[width=110mm]{gaze/app_E/7}
	\caption{ISE Generate Target PROM/ACE File画面}
%	\label{fig:　}
\end{figure}

\newpage

(8)~~ISE iMPACTが開く.

\begin{figure}[H]
	\centering
	\includegraphics[width=110mm]{gaze/app_E/8}
	\caption{ISE iMPACT画面}
%	\label{fig:　}
\end{figure}


(9)~~Create PROM Fileをダブルクリックする.

(10)~~PROM File Formatterが開く.以下の通りに選択する.\\
STEP1：Xilinx Flash/PROM\\
STEP2：PROM Family	→	Platform Flash\\
Dvice(bits)	→	xcf32p[32M]\\
ここでAdd Storage Deviceをクリックする.\\
STEP3：output File name→任意のファイル名(日本語が入らないようにすること)\\
output File Location→任意のフォルダを指定(日本語が入らないようにすること)\\
ここまで入力したら,OKをクリックする.\\
＊ここで作成される.mcsファイルのファイル名は,のちに選択するため覚えておく.\\
＊ファイル名は,日付と時間でつけると良い.(例：20200000$\_$0000)\\

\begin{figure}[H]
	\centering
	\includegraphics[width=110mm]{gaze/app_E/10}
	\caption{ISE Create PROM File画面}
%	\label{fig:　}
\end{figure}

\newpage

(11)~~Add DeviceはOKをクリックする.

\begin{figure}[H]
	\centering
	\includegraphics[width=110mm]{gaze/app_E/11}
	\caption{ISE Add Device画面}
%	\label{fig:　}
\end{figure}

(12)~~ .bitファイルを選択する.

\begin{figure}[H]
	\centering
	\includegraphics[width=110mm]{gaze/app_E/12}
	\caption{ISE Add Device選択画面}
%	\label{fig:　}
\end{figure}

\newpage

(13)~~Add DeviceダイアログはNOを選択する.そのあとのダイアログはOKを選択.

\begin{figure}[H]
	\centering
	\includegraphics[width=110mm]{gaze/app_E/13}
	\caption{ISE Add Device選択画面}
%	\label{fig:　}
\end{figure}

(14)~~ウィンドウ左側タブ内のGenerate Fileをダブルクリックする.

(15)~~Generate Succeededが表示される.

\begin{figure}[H]
	\centering
	\includegraphics[width=110mm]{gaze/app_E/15}
	\caption{ISE Generate File実行画面}
%	\label{fig:　}
\end{figure}


\newpage

(16)~~ウィンドウ左側タブ内のBoundary Scanをダブルクリックする.

\begin{figure}[H]
	\centering
	\includegraphics[width=110mm]{gaze/app_E/16}
	\caption{ISE Generate File実行画面}
%	\label{fig:　}
\end{figure}


(17)~~ウィンドウ右半分タブ内で右クリックし,Initialize Chainをクリックする.

\begin{figure}[H]
	\centering
	\includegraphics[width=110mm]{gaze/app_E/17}
	\caption{ISE Initialize Chain実行画面}
%	\label{fig:　}
\end{figure}


\newpage

(18)~~コントローラに入っているFPGAとPROMが検出されるのを確認したら,ダイアログのCancelをクリックする.

\begin{figure}[H]
	\centering
	\includegraphics[width=110mm]{gaze/app_E/18}
	\caption{ISE FPGA/PROM検出画面}
%	\label{fig:　}
\end{figure}

(19)~~xcf32pのブロックを右クリックし,Launch File Assignment Wizardをクリックする.

\begin{figure}[H]
	\centering
	\includegraphics[width=110mm]{gaze/app_E/19}
	\caption{ISE Launch File Assignment Wizard選択画面}
%	\label{fig:　}
\end{figure}

\newpage

(20)~~Assign New Configuration Fileが開くので,Bypassをクリックする.

\begin{figure}[H]
	\centering
	\includegraphics[width=110mm]{gaze/app_E/20}
	\caption{ISEAssign New Configuration File画面}
%	\label{fig:　}
\end{figure}


(21)~~手順(10)で作成した.mcsファイルを選択し,Openをクリックする.

\begin{figure}[H]
	\centering
	\includegraphics[width=110mm]{gaze/app_E/21}
	\caption{ISE .mcsファイル選択画面}
%	\label{fig:　}
\end{figure}

\newpage

(22)~~Device1はそのままで良い.Device2を選択し,追加で以下の項目にチェックを入れ,OKをクリックする.\\
Verify\\
Load FPGA\\
Parallel Mode\\

\begin{figure}[H]
	\centering
	\includegraphics[width=110mm]{gaze/app_E/22_1}
	\caption{ISE Device1画面}
%	\label{fig:　}
\end{figure}

\begin{figure}[H]
	\centering
	\includegraphics[width=110mm]{gaze/app_E/22_2}
	\caption{ISE Device2実行画面}
%	\label{fig:　}
\end{figure}

\newpage


(23)~~xcf32pのアイコンを右クリックし,programをクリックすると,書き込みが開始する.

\begin{figure}[H]
	\centering
	\includegraphics[width=110mm]{gaze/app_E/23}
	\caption{ISE 書き込み開始画面}
%	\label{fig:　}
\end{figure}


(24)~~ Program Succeededが表示されたら,書き込みが完了する.

\begin{figure}[H]
	\centering
	\includegraphics[width=110mm]{gaze/app_E/24}
	\caption{ISE 書き込み完了画面}
%	\label{fig:　}
\end{figure}

\newpage

(25)~~コントローラボードのLED6とLED8が点灯したら書き込み成功となる.\\
＊点灯していなかった場合,書き込みが出来ていない可能性があるので,もう一度やってみるか,プログラムを見直す(仕様メモリを減らす,行数をなるべく減らす等する)必要がある.

(26)~~24V電源を切り,JTAGを外す.\\
*DSPの接続をしたまま書き込みをしているなら,そのまま実機検証できる.外す必要はない.\\
基本的に,コントローラ上の配線を抜き差しする場合は,電源を切ること.


\newpage
\section{ロジックアナライザの使用方法}
(1)~~Analyze Design Using Chipscopeをダブルクリックする．

\begin{figure}[H]
	\centering
	\includegraphics[width=110mm]{gaze/app_E/chipscope1}
	\caption{chipscope}
%	\label{fig:　}
\end{figure}

(2)~~左上にあるOpen cable /Search JTAG chain を押し初期化を行う．ここでJTAG エラーと表示される場合がある．その時はSDK 上でXilinx Tools ＞ Configure JTAG settings を選択し，JTAG type をDigilent USB Cable に変更することで回避できる．

\begin{figure}[H]
	\centering
	\includegraphics[width=110mm]{gaze/app_E/chipscope2}
	\caption{chipscope}
%	\label{fig:　}
\end{figure}

(3)~~Device → configureを選択し，bitファイルを読み込む．

\begin{figure}[H]
	\centering
	\includegraphics[width=110mm]{gaze/app_E/chipscope3}
	\caption{chipscope}
%	\label{fig:　}
\end{figure}

(4)~~Trigger Setup，Waveform，Bus Plotをダブルクリックすると，それぞれトリガ設定，内部信号の数値，波形を見ることができる．

\begin{figure}[H]
	\centering
	\includegraphics[width=110mm]{gaze/app_E/chipscope4}
	\caption{chipscope}
%	\label{fig:　}
\end{figure}

(5)~~Trigger Setupでは測定信号のトリガを設定できる．

\begin{figure}[H]
	\centering
	\includegraphics[width=110mm]{gaze/app_E/chipscope5}
	\caption{chipscope}
%	\label{fig:　}
\end{figure}

(6)~~Waveformでは取得した信号の数値を見ることができ，符号の有無等の数値の表示形式を変更できる．

\begin{figure}[H]
	\centering
	\includegraphics[width=110mm]{gaze/app_E/chipscope6}
	\caption{chipscope}
%	\label{fig:　}
\end{figure}

(7)~~波形の取得方式として，単発取得，連続取得を決めることができる．

\begin{figure}[H]
	\centering
	\includegraphics[width=110mm]{gaze/app_E/chipscope8}
	\caption{chipscope}
%	\label{fig:　}
\end{figure}

\begin{figure}[H]
	\centering
	\includegraphics[width=110mm]{gaze/app_E/chipscope9}
	\caption{chipscope}
%	\label{fig:　}
\end{figure}

\begin{comment}


\end{comment}

